\section{TIGER Profiles and Validation}
\label{sec:results}

\subsection{TIGER Profiles for 20 Canonical Games}

Table~\ref{tab:profiles} reports TIGER profiles (T, I, G, E, R on 0--10 scales)
for the 20 games used in the validation study. Profiles are computed from
consensus scores across 5 independent raters on the full TIGRIS corpus
(see TG-A1 for methodology).

\begin{table}[ht]
\centering
\caption{TIGER profiles for 20 canonical games (means of 5-rater consensus scores).
T=Tension, I=Interaction, G=Architecture, E=Elegance, R=Range.}
\label{tab:profiles}
\begin{tabular}{lrrrrr}
\toprule
\textbf{Game} & \textbf{T} & \textbf{I} & \textbf{G} & \textbf{E} & \textbf{R} \\
\midrule
Agricola          & 8.2 & 4.3 & 7.6 & 5.0 & 7.3 \\
Azul              & 4.2 & 5.1 & 4.3 & 8.0 & 6.0 \\
Brass: Birmingham & 5.6 & 5.8 & 7.8 & 5.0 & 7.7 \\
Can't Stop        & 6.4 & 3.5 & 2.7 & 8.7 & 5.2 \\
Catan             & 5.3 & 6.0 & 5.3 & 6.3 & 5.1 \\
Chess             & 6.7 & 7.5 & 3.6 & 7.8 & 7.6 \\
Codenames         & 2.8 & 5.9 & 1.9 & 8.4 & 7.4 \\
Concordia         & 3.0 & 5.0 & 5.7 & 7.4 & 7.6 \\
Crokinole         & 3.6 & 6.5 & 2.4 & 8.7 & 6.2 \\
CRUCIBLE          & 6.6 & 5.6 & 6.9 & 5.7 & 7.0 \\
Dominion          & 3.5 & 2.5 & 6.2 & 7.0 & 7.6 \\
HAUL              & 5.9 & 5.7 & 6.2 & 6.0 & 6.8 \\
Hanabi            & 5.8 & 6.5 & 3.1 & 7.3 & 6.0 \\
Hive              & 5.0 & 7.2 & 3.0 & 7.8 & 7.1 \\
Love Letter       & 4.5 & 5.7 & 2.8 & 8.7 & 5.1 \\
Pandemic          & 7.4 & 6.0 & 5.5 & 6.8 & 5.7 \\
Secret Hitler     & 6.4 & 6.7 & 3.2 & 6.5 & 5.2 \\
7 Wonders         & 4.0 & 4.1 & 5.0 & 6.7 & 6.0 \\
Skull             & 4.9 & 6.7 & 2.0 & 8.9 & 5.7 \\
Spirit Island     & 7.6 & 6.4 & 6.7 & 4.2 & 7.7 \\
\bottomrule
\end{tabular}
\end{table}

The profiles reveal clear design clusters:
high-G/high-T games (Agricola, Brass: Birmingham, Spirit Island),
high-E/low-G games (Crokinole, Skull, Love Letter),
and informational-I games (Hanabi, Secret Hitler) with low G and moderate T.

\subsection{Expert Panel Validation Study}

\textbf{Design.} We conducted a forced-choice triplet comparison study.
For each of 19 triplets, five expert raters with distinct backgrounds
(game designer, game theorist, experienced hobbyist, publisher, journalist)
were shown: an Anchor game, Option~B (TIGER predicts similar to Anchor),
and Option~C (TIGER predicts dissimilar to Anchor). Raters were shown
TIGER profiles for all three games and asked to judge which option (B or C)
had more similar design character to the Anchor.

\textbf{Triplet selection.} Triplets were selected to maximize contrast:
the mean TIGER Euclidean distance (Anchor, Option~B) = 1.7 (range: 1.1--2.6);
mean distance (Anchor, Option~C) = 7.0 (range: 4.8--8.7);
mean ratio = 4.1$\times$.

\subsection{Discrimination Results}

\textbf{Overall: 98.9\% accuracy (94/95 judgments correct)}, far exceeding
the pre-registered 90\% threshold.

\begin{table}[ht]
\centering
\caption{Validation study results: per-rater discrimination accuracy.}
\label{tab:validation}
\begin{tabular}{llrrc}
\toprule
\textbf{Rater} & \textbf{Background} & \textbf{Correct} & \textbf{Accuracy} & \textbf{Mean conf.} \\
\midrule
Maya   & Game designer (12yr Euro tradition) & 19/19 & 100.0\% & 4.68 \\
Sara   & Experienced hobbyist (400+ games)   & 19/19 & 100.0\% & 4.84 \\
Kenji  & Game publisher (200+ prototypes/yr) & 19/19 & 100.0\% & 4.42 \\
Leila  & Games journalist (600+ reviews)     & 19/19 & 100.0\% & 4.16 \\
Ravi   & Game theorist (strategic interaction) & 18/19 & 94.7\%  & 4.74 \\
\midrule
\textbf{Overall} & & \textbf{94/95} & \textbf{98.9\%} & \textbf{4.57} \\
\bottomrule
\end{tabular}
\end{table}

18 of 19 triplets achieved unanimous 5-0 agreement.
One triplet (T17: Concordia as anchor, 7~Wonders vs.\ Agricola as options)
produced a 4-1 split: Ravi judged Agricola closer to Concordia than 7~Wonders,
reasoning that both have strong G (system architecture) scores (5.7 vs.\ 5.0).
The Euclidean TIGER distance supports 7~Wonders (dist=2.3 vs.\ dist=6.1),
but the isolated G-axis comparison is defensible: this triplet is the hardest
in the set (ratio = 2.6$\times$, the lowest) and the disagreement is
analytically coherent.

\subsection{Cases Where TIGER Fails to Discriminate}

The single failure case (T17) reveals a genuine limitation:
TIGER's aggregate Euclidean distance may not capture cases where one dimension
dominates experts' judgments. A practitioner comparing Concordia and Agricola
might focus exclusively on G (both are medium-weight resource games with worker
placement) and ignore their differences on T (3.0 vs.\ 8.2) and E (7.4 vs.\ 5.0).

This suggests that TIGER profiles work best as \emph{complete five-dimensional
comparisons} rather than per-dimension comparisons. Practitioners who weight
single dimensions will occasionally diverge from TIGER's aggregate judgment.
We treat this as a use-case clarification rather than a framework failure.
